\chapter{Validité scientifique}

\section{Objet de l’analyse}

Les résultats précédents couvrent des objets différents : performance prédictive, reconnaissance de format, exécution distribuée et corrélation. Leur validité ne peut pas être évaluée par un score unique. Ce chapitre examine les menaces internes, externes, de construit et statistiques, puis précise les conclusions qui résistent à ces contrôles.

\section{Validité interne}

La validité interne concerne le lien entre le protocole et le résultat observé. Le principal risque est une contamination entre entraînement et test. Sur HDFS, les blocs sont disjoints ; Drain3 apprend 28 clusters sur l’entraînement, puis utilise \texttt{match\_only}. Le hash de l’état reste identique avant et après l’inférence. L’agrégateur et le seuil sont sélectionnés sur validation. Ces assertions soutiennent l’absence des fuites précisément instrumentées.

CICIDS2017 illustre un autre risque : un split aléatoire place des observations de même scénario dans les deux ensembles. Le score de 0,995142 reste correct pour cette question statistique, mais il ne peut pas servir de preuve de transfert. Les holdouts scénario et temporel sont les contrôles pertinents, avec des F1 nettement plus faibles.

Sur CSE-CIC-IDS2018, le split par jour limite le mélange direct des lignes. Le score mono-feature de \texttt{Dst Port} et la chute après ablation montrent cependant une dépendance structurelle. Coefficients, permutation et retraits successifs n’identifient pas de fuite triviale correspondant aux mécanismes testés. Ils n’excluent pas un artefact propre aux journées ou à la génération du dataset.

Le routeur open-set est calibré sur les sources connues seulement. Les trois fichiers inconnus ne servent pas au choix du seuil 100. Cette séparation protège l’évaluation immédiate, mais la taille du jeu inconnu reste trop faible pour caractériser toute nouveauté possible.

\section{Validité externe}

Les datasets publics ne reproduisent pas la diversité d’un réseau d’entreprise contemporain. Les résultats CICIDS2017 et CSE-CIC-IDS2018 dépendent de scénarios construits et de caractéristiques de flux déjà calculées. Les scores ne se transfèrent donc pas automatiquement à d’autres capteurs, topologies ou périodes.

HDFS et BGL représentent des systèmes particuliers. Le résultat HDFS porte sur 2,000 blocs de test et 29 positifs. Le bootstrap décrit la sensibilité à ce test, non une population universelle. Sur BGL, 89,8081\,\% des unités de test comportent un template inconnu ; un autre taux de nouveauté pourrait modifier fortement le résultat.

La campagne multi-VM utilise deux systèmes Linux et un Redis central pendant 2,712455~s. Elle établit une distribution inter-machine et une reprise contrôlée. Elle ne couvre pas la haute disponibilité, les partitions, le multi-site, le vieillissement des agents ou une charge prolongée. Les conclusions restent volontairement locales.

\section{Validité de construit}

Une métrique n’est valable que si elle représente le concept annoncé. Sur HDFS, le label porte sur le bloc : le F1 événementiel n’est pas l’indicateur principal. Pour BGL, un F1 global sans groupe connu/inconnu confondrait anomalie et nouveauté de template. Pour le routeur, la \texttt{decision\_margin} est une marge heuristique, pas une confiance probabiliste.

La couverture multi-format mesure trois opérations : lire, parser et normaliser les unités présentes. Elle ne mesure ni la complétude de chaque champ, ni la fidélité bit à bit au journal brut, ni la compatibilité avec toutes les implémentations d’un format. La fixture Apache d’une ligne doit rester visible dans l’interprétation.

Le caractère multi-agent est évalué par des comportements : capacités, état, plusieurs handlers, propositions, refus, attributions et résultats. Le nombre de processus seul ne serait pas une preuve suffisante. De même, un replay idempotent protège un effet logique ; il ne démontre pas une transaction distribuée générale.

Enfin, l’exécution de modèles réels et la qualité de leurs prédictions sont deux construits distincts. La campagne M démontre 1,400 inférences effectives sur 1,401 unités, mais aucune exactitude globale n’est calculée en l’absence de labels homogènes.

\section{Validité statistique}

Les campagnes CICIDS2017 et CSE-CIC-IDS2018 utilisent cinq graines, tandis que le benchmark A--D répète dix fois chaque combinaison. Ces répétitions documentent la dispersion, mais cinq graines restent insuffisantes pour décrire toutes les variations de scénario. L’écart-type élevé des holdouts CICIDS2017 doit être lu comme une hétérogénéité des attaques tenues hors apprentissage.

Le bootstrap HDFS compte 1,000 réplications au niveau bloc, unité indépendante du test. Son intervalle empirique reflète le faible nombre de positifs. Il ne corrige pas un éventuel biais de sélection du corpus. BGL ne permet aucun rappel pour le groupe connu, car celui-ci ne contient aucun positif ; inscrire zéro comme rappel observé sans cette précision serait trompeur.

Le benchmark multi-agent mesure des micro-tâches. Le monolithe bénéficie d’un coût de coordination minimal, alors que le Contract Net écrit plusieurs messages par attribution. Une charge plus lourde pourrait déplacer le compromis, mais cette possibilité n’est pas un résultat. Les données présentes soutiennent seulement l’absence de gain systématique de C ou D sur le protocole exécuté.

\section{Risques résiduels et mesures de contrôle}

\begin{table}[H]
\centering
\caption{Principaux risques résiduels après consolidation.}
\label{tab:risques-validite}
\small
\begin{tabularx}{\textwidth}{p{0.22\textwidth} p{0.31\textwidth} X}
\toprule
\textbf{Risque} & \textbf{Contrôle réalisé} & \textbf{Limite restante} \\
\midrule
Proximité train/test & Splits scénario, temporel et par bloc & Datasets publics et scénarios limités. \\
Parsing contaminé & Drain3 train-only, état gelé & Vérifié pour le pipeline HDFS final. \\
Nouveauté lexicale BGL & Groupes connu/inconnu et baseline & Aucun positif dans le groupe connu. \\
Dépendance aux variables & Mono-features, coefficients, ablation, permutation & Biais journalier toujours possible. \\
Surajustement du rejet & Calibration sur connus seulement & Trois inconnus seulement. \\
Confusion distribution/HA & Deux VM et Redis central explicités & Broker et partitions non testés. \\
Confusion mémoire/gain & Comparaison ON/OFF & Aucun bénéfice systématique observé. \\
\bottomrule
\end{tabularx}
\end{table}

\section{Conclusions soutenues}

Les preuves permettent d’affirmer que le prototype normalise les unités des sources testées, route correctement les 28 fichiers connus du jeu indépendant, rejette trois inconnus au seuil gelé et exécute les modèles réels pour 1,400 unités sur 1,401. Elles permettent aussi d’affirmer que les agents proposent, refusent, exécutent et réattribuent des tâches sur deux VM, avec un replay sans effet dupliqué.

Elles ne permettent pas d’annoncer une exactitude multi-source, une supériorité générale du routage spécialisé, un gain systématique de la mémoire, une mise à jour autonome évaluée ou une disponibilité industrielle. Cette frontière n’est pas un défaut rédactionnel : elle constitue le résultat de l’audit scientifique.
